\chapter{Introduction}\label{ch:intro}

\textsf{KELP} is a Lagrangian particle-tracking engine for the velocity
fields produced by unstructured-mesh ocean circulation models, with first-class
support for the FESOM family: the Finite-volumE Sea ice--Ocean Model in its second
version \citep[FESOM2;][]{danilov2017fesom2}, on purely triangular meshes, and its
coastal relative FESOM-C, on mixed triangular--quadrilateral meshes. It also reads ocean fields
remapped to the HEALPix grid and, in a first two-dimensional version, structured NEMO grids. Given a time sequence of stored model
velocity snapshots and a set of release positions, it integrates the trajectories
of passive particles either forward or backward in time, optionally in three
dimensions and optionally subject to a stochastic eddy-diffusion term. This manual
documents the numerical methods that the engine implements and links each method
to the Python source that realizes it.

\section{The Lagrangian transport problem}\label{sec:intro-lagrangian}

In the Lagrangian description of transport one follows individual fluid parcels
along their paths, in contrast to the Eulerian description, which records the flow
at fixed locations. A passive particle carried by a velocity field
$\velvec(\pos,t)$ obeys the ordinary differential equation
\begin{equation}\label{eq:lagrangian-ode}
  \frac{\dd \pos}{\dd t} = \velvec\!\bigl(\pos(t),\,t\bigr),
  \qquad \pos(t_0)=\pos_0 ,
\end{equation}
where the position $\pos=(\lonx,\latx,\depx)$ collects longitude, latitude and
depth, and the velocity $\velvec=(u,v,\Wvel)$ collects the two horizontal
components and the vertical component. We adopt throughout the convention that
depth $\depx$ and vertical velocity $\Wvel$ are positive downward, so that a
positive $\Wvel$ deepens a particle. The particle is passive: it does not modify
the flow, and \cref{eq:lagrangian-ode} is integrated as a one-way coupling from
the stored model fields to the trajectory.

Solving \cref{eq:lagrangian-ode} numerically poses two distinct subproblems.
First, the right-hand side must be evaluated at the particle position $\pos(t)$,
which in general falls between the mesh vertices at which the model stores its
fields and between the discrete times at which snapshots are available; this
requires interpolation in space and in time. Second, the differential equation
must be advanced from one time level to the next with a controlled truncation
error; this requires a time-integration scheme. The spatial-interpolation
subproblem is the subject of \cref{ch:hvel,ch:vertical,ch:sampling}, the
temporal-reconstruction subproblem of \cref{ch:temporal}, and the
time-integration subproblem of \cref{ch:integration}. When the stochastic term is included, a
backward run integrates the adjoint process --- the advection reversed, the diffusion not --- so that
backward integration remains a well-posed random walk (\cref{ch:diffusion}). The remaining chapters treat
the supporting machinery: locating a particle on the mesh (\cref{ch:localization}),
the stochastic-diffusion extension of \cref{eq:lagrangian-ode}
(\cref{ch:diffusion}), the treatment of domain boundaries and the sea floor
(\cref{ch:lifecycle}), and the reading, writing, configuration and verification of
the whole (\crefrange{ch:input}{ch:performance}); a final part describes the agent skills through which
runs are set up and analysed (\cref{ch:skills}).

For a broad survey of Lagrangian ocean analysis --- its applications, its
numerical pitfalls, and the conventions adopted in the field --- we refer the
reader to \citet{vansebille2018lagrangian}; the present manual is concerned
specifically with how \textsf{KELP} solves \cref{eq:lagrangian-ode} on an
unstructured mesh.

\section{Unstructured-mesh ocean models and the FESOM family}\label{sec:intro-fesom2}

FESOM2 discretizes the ocean horizontally on an unstructured triangular mesh and
vertically on a set of fixed depth levels \citep{danilov2017fesom2}; its coastal
relative FESOM-C uses a \emph{mixed} mesh, in which an element may be a triangle or
a quadrilateral, together with a terrain-following vertical coordinate.
\textsf{KELP} supports both, and represents the other sources on the same kind of mesh: a HEALPix
product becomes either a mesh whose nodes are the pixel centres (quadrilaterals, with triangles at the
eight points where only three pixels meet) or a mesh whose elements are the pixels themselves, and a
structured NEMO grid becomes a mesh of quadrilaterals (\cref{ch:mesh}). We therefore treat an element as
a triangle or a quadrilateral throughout. The unstructured
horizontal mesh permits the spatial resolution to vary smoothly across the domain
--- fine where the dynamics demand it, coarse elsewhere --- but it removes the
regular row-and-column addressing of a latitude--longitude grid. As a result,
identifying the element that contains a given particle position is itself a
search problem rather than an arithmetic indexing operation; we treat it in
\cref{ch:localization}.

A second feature of the unstructured mesh bears directly on interpolation: the
model stores different quantities at different geometric locations. The two
horizontal velocity components may be defined either at the element vertices
(``node-layout'' fields, such as the FESOM2 variables \conf{unod} and
\conf{vnod}, the HEALPix pixel-centre mesh and the NEMO grids) or at the element centroids
(``element-layout'' fields, such as the FESOM2 variables \conf{u} and \conf{v} and the FESOM-C
velocities). The
vertical velocity and other interface quantities are defined at the vertical
interfaces that bound the depth layers, while scalar fields such as temperature
and salinity are defined at the nodes. Each storage location calls for a
different interpolation stencil, a distinction that recurs throughout
\crefrange{ch:hvel}{ch:sampling} and that we render explicitly in every schematic.

The engine ships with ready configurations that make these choices concrete: global two-dimensional
runs on node velocities, a regional three-dimensional run with vertical velocity and sampled temperature
and salinity, a structured-grid run and HEALPix runs. Where this manual quotes a size it is usually that of
the FESOM2 \emph{core2} global mesh (\num{126858} nodes, 48 vertical interfaces reaching to about
\SI{6250}{\metre}). The object of this document is the engine, not any experiment run with it.

\section{Scope and purpose}\label{sec:intro-scope}

This manual is a numerical and implementation reference. For each component we
state the governing equations, give the discrete formulae exactly as implemented
(including the sign conventions, the clamping rules, and the degenerate-case
fallbacks that a working code must contain), illustrate the geometry of the
computation with a schematic, and cite the source location. We do not report
scientific results, and we do not attempt to assess the physical realism of any
particular simulation; our goal is to characterize what the engine computes and
how, not to interpret the trajectories it produces.

The engine is accompanied by an extensive automated test suite of about \num{2500} tests. The large majority of these guard low-level invariants --- array
shapes, configuration parsing, byte-for-byte reproducibility --- and are not
individually informative. A small number, by contrast, verify the numerical
physics against independent references: convergence under step refinement in
time-dependent eddy flows, conservation of a
vertical streamfunction under continuity-consistent flow, reversibility of
the integrator under time reversal (to within a small tolerance), and agreement with closed-form analytic flows.
We discuss these by name in \cref{ch:validation} as evidence that the numerics
behave as specified, without reproducing their internals.

\section{How this manual is organized}\label{sec:intro-org}

The chapters follow the order in which the engine processes a single integration
step, so that reading the manual front to back traces one time step through the
solver. \Cref{ch:arch} gives the overall architecture --- the data structures and
the per-step pipeline --- and serves as the map to which the later chapters
refer. The pairs and groups that follow treat, in turn, the mesh and particle
localization; spatial interpolation of velocity and of scalars; time integration
and temporal reconstruction; stochastic diffusion and the particle lifecycle; and
finally the data, configuration, validation and performance layers. A last part, ``Working with
KELP'', describes the agent skills that turn a described experiment into a run and its output into
figures and statistics (\cref{ch:skills}). The appendices provide an exhaustive configuration-key reference, the symbol list,
the source-code map, and a glossary. Readers seeking a particular topic should
consult the front-matter guide ``How to use this manual'' or the source-code map
in \cref{app:codemap}.

\section{Conventions}\label{sec:intro-conv}

We summarize the principal conventions here and state them in full in the
front-matter section ``Notation and conventions''. Depth and vertical velocity are
positive downward. Angles are in degrees and velocities in SI units unless stated
otherwise. The integration direction is carried by the factor $\tdir=\pm1$, with
the time step $\dt$ kept positive. The spherical metric converts advective velocities to
angular increments through $\zzero=\cos(\latx)\,\Mlat$ and
$\Mlat=\SI{111120}{\metre\per\degree}$; diffusive kicks move along great circles instead
(\cref{sec:diff-sphere}). Model files that store the vertical velocity positive upward are converted when
they are read (\cref{sec:vert-wsign}). Fields are identified throughout as living
at nodes, at element centres, or at vertical interfaces. The complete symbol list
is given in \cref{app:notation}.
